%% kapa: exam=1 round=23 year=2012 period=2 form=A subject=appraisal_law subject_ko=부동산관계법규 seq=28 number=68
%% source: Q-Net 감정평가사 기출문제 (cst00309) artlSeq=5052580 file=raw/1008/5052580/2012년 제23회 감정평가사 제1차시험 기출문제(2교시 A형).pdf
%% conversion: kordoc 4.24.51 -> tools/gen1.py (auto); review=proofread-sonnet; flags=none
\begin{question}{68}
국유재산법령상 행정재산의 사용허가 등에 관한 설명으로 옳은 것은?
\begin{choices}
  \item 공유(公有) 또는 사유재산과 교환하여 그 교환받은 재산을 행정재산으로 관리하려는 경우에는 행정재산도 교환하거나 양여할 수 있다.
  \item 행정재산의 사용허가에 관하여는「국가를 당사자로 하는 계약에 관한 법률」의 규정이 준용되지 않는다.
  \item 행정재산을 주거용으로 사용허가할 수 있지만, 수의의 방법으로 할 수는 없다.
  \item 5년을 초과하지 아니하는 범위에서 행정재산의 사용허가를 갱신할 수 있지만, 수의의 방법으로 사용허가를 받은 재산에 대하여는 1회에 한하여 갱신할 수 있다.
  \item 중앙관서의 장이 행정재산을 직접 공용으로 사용하려는 지방자치단체에게 사용허가하는 경우에는 사용료를 면제하여야 한다.
\end{choices}
\end{question}
\answerall
\answernote{정답표에 1,2,3,4,5 모두 정답으로 표기(전원 정답)}
